
% ---------------------------------------------------------
% Back cover
% ---------------------------------------------------------

% Finish the preceding content.
\clearpage

% The back cover must be on an even (left-hand) page.
% If the next page is odd, insert a blank page.
\ifodd\value{page}
  \thispagestyle{empty}
  \null
  \clearpage
\fi

% ---------------------------------------------------------
% Back cover page
% ---------------------------------------------------------

\thispagestyle{empty}

\vspace*{15mm}


{
\fontsize{18pt}{22pt}\selectfont \bfseries\sffamily
Über dieses Buch
\par
}

\vspace{6mm}

{\fontsize{11pt}{16pt}\selectfont\bfseries\sffamily Problemstellung:\par}
{\fontsize{11pt}{16pt}\selectfont\sffamily
Magnetische Maßverkörperungen bilden die Grundlage zahlreicher berührungsloser Weg- und Winkelmesssysteme. An ihrer Entwicklung, Herstellung, Integration und Prüfung sind meist mehrere Akteure entlang der Wertschöpfungskette beteiligt. Einheitliche Begriffe und Darstellungsstandards sind entscheidend für eine klare technische Kommunikation, verringern den Entwicklungsaufwand und ermöglichen eine vergleichbare Qualitätsbewertung.
\par
}

\vspace{3mm}

{\fontsize{11pt}{16pt}\selectfont\bfseries\sffamily Inhalt:\par}
{\fontsize{11pt}{16pt}\selectfont\sffamily
Diese Richtlinie schafft einen gemeinsamen Rahmen für die technische Beschreibung magnetischer Maßverkörperungen. Sie greift etablierte Fachbegriffe auf, vereinheitlicht deren Verwendung und führt physikalisch formalisierte und in sich konsistente Definitionen ein. Ein besonderer Schwerpunkt liegt auf der Ausarbeitung praxisrelevanter mechanischer und magnetischer Bezugssysteme, die als Grundlage für die technische Spezifikation und Vermessung dienen. Darauf aufbauend definiert die Richtlinie Konventionen zur Darstellung magnetischer Strukturen und Streufelder sowie deren charakteristische Merkmale und Kenngrößen. Die vorgeschlagene systematische Einteilung und Kurznotation gängiger Maßverkörperungen erleichtert die Kommunikation. Anschauliche Abbildungen und Anwendungsbeispiele unterstützen die Erläuterungen.
\par
}

\vspace{3mm}

{\fontsize{11pt}{16pt}\selectfont\bfseries\sffamily Zielgruppen:\par}
{\fontsize{11pt}{16pt}\selectfont\sffamily
Die Richtlinie richtet sich an Hersteller, Entwickler, Prüfeinrichtungen und Anwender magnetischer Messsysteme. Darüber hinaus bietet sie Technikern, Studierenden und Einsteigern einen verständlichen Zugang zu den physikalischen Grundlagen und technischen Konzepten magnetischer Maßverkörperungen in Weg- und Winkelmesssystemen.
\par
}

\vfill

{
\sffamily\textbf{INNOMAG e. V.}\\
https://innomag.org/
}

\vspace*{10mm}


% ---------------------------------------------------------
% Backcover background: horizontal boundary
% ---------------------------------------------------------

\AddToShipoutPictureBG*{%
  \begin{tikzpicture}[remember picture,overlay]

    % Upper area
    \fill[titlebg1]
      (current page.south west)
      rectangle
      (current page.north east);

    % Lower area: horizontal continuation
    \fill[titlebg2]
      (current page.south west)
      --
      ([yshift=-0.6\paperheight]current page.north west)
      --
      ([yshift=-0.6\paperheight]current page.north east)
      --
      (current page.south east)
      -- cycle;

  \end{tikzpicture}%
}